% Manual Técnico — Planck Automation
% Compilar com: pdflatex manual_tecnico.tex  (rodar duas vezes pelo sumário)
\documentclass[11pt,a4paper]{article}

\usepackage[utf8]{inputenc}
\usepackage[T1]{fontenc}
\usepackage[brazil]{babel}
\usepackage[margin=2.5cm]{geometry}
\usepackage{amsmath}
\usepackage{amssymb}
\usepackage{booktabs}
\usepackage{xcolor}
\usepackage{enumitem}
\usepackage{longtable}
\usepackage{listings}
\usepackage[colorlinks=true,linkcolor=azulesc,urlcolor=azulesc]{hyperref}

\definecolor{azulesc}{HTML}{0D47A1}
\definecolor{fundocod}{HTML}{F5F5F5}
\definecolor{comentario}{HTML}{558B2F}
\definecolor{palavra}{HTML}{1565C0}

\lstset{
  basicstyle=\ttfamily\footnotesize,
  backgroundcolor=\color{fundocod},
  keywordstyle=\color{palavra}\bfseries,
  commentstyle=\color{comentario}\itshape,
  stringstyle=\color{black},
  showstringspaces=false,
  breaklines=true,
  frame=leftline,
  framesep=6pt,
  xleftmargin=10pt,
  literate={á}{{\'a}}1 {é}{{\'e}}1 {í}{{\'i}}1 {ó}{{\'o}}1 {ú}{{\'u}}1
           {ã}{{\~a}}1 {õ}{{\~o}}1 {ç}{{\c c}}1 {â}{{\^a}}1 {ê}{{\^e}}1
           {Á}{{\'A}}1 {É}{{\'E}}1 {ô}{{\^o}}1 {à}{{\`a}}1
}

\newcommand{\arqv}[1]{\texttt{#1}}
\newcommand{\classe}[1]{\texttt{#1}}

\newenvironment{nota}
  {\par\medskip\noindent\begin{tabular}{|p{0.96\linewidth}|}\hline\rule{0pt}{3ex}}
  {\\[1ex]\hline\end{tabular}\par\medskip}

\title{\textbf{Planck Automation}\\[0.3em]
\large Determinação da constante de Planck por radiação de corpo negro\\[0.6em]
\Large Manual Técnico}
\author{Arquitetura, modelo físico e lógica interna}
\date{\today}

\begin{document}
\maketitle
\thispagestyle{empty}

\begin{abstract}
\noindent
Este documento descreve como o software funciona por dentro: a arquitetura em
camadas, o modelo de concorrência, a cadeia de cálculo que converte medidas
elétricas na constante de Planck, os drivers dos instrumentos e a bancada
simulada. Destina-se a quem vai ler, manter ou estender o código. Para operar o
software, consulte o \textbf{Manual do Usuário}.
\end{abstract}

\tableofcontents
\newpage

% =============================================================
\section{Arquitetura}

\subsection{Camadas}

O software se organiza em quatro camadas, com dependências apontando sempre
para baixo:

\begin{center}
\begin{tabular}{@{}p{3.0cm}p{11.5cm}@{}}
\toprule
\textbf{Camada} & \textbf{Responsabilidade e conteúdo} \\
\midrule
\arqv{ui/} &
Tudo que o usuário vê: janela principal, abas, painéis, tema. Não conhece
VISA nem faz contas de física. \\[0.4em]
\arqv{core/} &
Comunicação com os instrumentos: drivers SCPI, varredura e validação de
recursos, bancada simulada. \\[0.4em]
\arqv{utils/} &
Funções puras: o modelo físico-matemático e a geração do relatório PDF. Não
dependem de Qt nem de hardware. \\[0.4em]
\arqv{content/} &
Dados tipados que não são código --- hoje, a lista de referências
bibliográficas. \\
\bottomrule
\end{tabular}
\end{center}

\subsection{Mapa de arquivos}

\begin{center}
\begin{longtable}{@{}p{5.6cm}r p{7.2cm}@{}}
\toprule
\textbf{Arquivo} & \textbf{Linhas} & \textbf{Papel} \\
\midrule
\endhead
\arqv{main.py} & 18 & Ponto de entrada: cria o \classe{QApplication}, aplica o estilo Fusion e o tema escuro, abre a janela. \\[0.3em]
\arqv{ui/main\_window.py} & 59 & \classe{MainWindow}: monta as quatro abas e garante o desligamento seguro ao fechar. \\[0.3em]
\arqv{ui/theme.py} & 119 & Folha de estilos QSS do tema escuro. \\[0.3em]
\arqv{ui/components/} \arqv{connection\_panel.py} & 171 & Painel de ligações: varredura, validação, modo demonstração e persistência das preferências. \\[0.3em]
\arqv{ui/components/} \arqv{export\_dialog.py} & 39 & Janela de metadados do relatório. \\[0.3em]
\arqv{ui/tabs/tab\_simulation.py} & 325 & \classe{SimulationWorker} e a interface da simulação analítica. \\[0.3em]
\arqv{ui/tabs/tab\_experiment.py} & 445 & \classe{ExperimentWorker} e a interface da coleta na bancada. \\[0.3em]
\arqv{ui/tabs/tab\_references.py} & 127 & Renderiza as fichas de referência. \\[0.3em]
\arqv{core/hardware\_manager.py} & 257 & Drivers SCPI reais, threads de varredura e validação, seleção de drivers. \\[0.3em]
\arqv{core/mock\_hardware.py} & 300 & Bancada simulada e drivers falsos. \\[0.3em]
\arqv{utils/math\_models.py} & 86 & O modelo físico: temperatura, simulação e regressão. \\[0.3em]
\arqv{utils/pdf\_exporter.py} & 79 & Montagem do relatório PDF com reportlab. \\[0.3em]
\arqv{content/referencias.py} & 145 & Lista tipada das referências. \\
\bottomrule
\end{longtable}
\end{center}

\subsection{Ciclo de vida}

\arqv{main.py} faz três coisas e sai do caminho:

\begin{lstlisting}[language=Python]
app = QApplication(sys.argv)
app.setStyle("Fusion")          # motor de renderização neutro
app.setStyleSheet(DARK_THEME)   # QSS do tema escuro
window = MainWindow()
window.show()
sys.exit(app.exec())
\end{lstlisting}

\classe{MainWindow} instancia um único \classe{HardwareManager} e o compartilha
entre o painel de ligações e a aba de experimento. Ao fechar a janela, o
\texttt{closeEvent} verifica se há coleta em curso; se houver, sinaliza a
parada e \emph{aguarda} a thread terminar, para que a fonte seja desligada
antes de o processo morrer.

% =============================================================
\section{Configuração persistente}

O software guarda preferências no registro do Windows via \classe{QSettings},
sob a organização \texttt{Senac} e a aplicação \texttt{PlanckAutomation}. As
chaves usadas:

\begin{center}
\begin{tabular}{@{}p{5.4cm}p{9.2cm}@{}}
\toprule
\textbf{Chave} & \textbf{Conteúdo} \\
\midrule
\texttt{Connection/LastPWSName} & Nome amigável da última fonte validada \\
\texttt{Connection/LastPWSRes} & Endereço VISA da última fonte validada \\
\texttt{Connection/LastDMMName} & Nome amigável do último multímetro validado \\
\texttt{Connection/LastDMMRes} & Endereço VISA do último multímetro validado \\
\texttt{Connection/DemoMode} & Booleano: modo demonstração ligado \\
\bottomrule
\end{tabular}
\end{center}

Os nomes da organização e da aplicação estão centralizados em
\arqv{core/hardware\_manager.py} (constantes \texttt{ORGANIZACAO} e
\texttt{APLICACAO}), acessíveis pela função \texttt{preferencias()}.

\begin{nota}
\textbf{Acoplamento a observar.} A aba de experimento lê os endereços VISA
diretamente do \classe{QSettings}, e não por sinal vindo do painel de ligações.
É um canal lateral: funciona, mas cria uma dependência invisível entre os dois
componentes. Endereços de instrumentos simulados nunca são gravados nessas
chaves, para que desligar o modo demonstração devolva o último instrumento
real.
\end{nota}

% =============================================================
\section{Camada de hardware}

\subsection{O driver da fonte}

\classe{PWS4323\_Driver} encapsula os comandos SCPI da fonte:

\begin{center}
\begin{tabular}{@{}llp{6.6cm}@{}}
\toprule
\textbf{Método} & \textbf{Comando SCPI} & \textbf{Efeito} \\
\midrule
\texttt{configure\_safety\_limits} & \texttt{SOUR:CURR} & Define o limite de corrente \\
\texttt{set\_output} & \texttt{OUTP ON/OFF} & Liga ou desliga a saída \\
\texttt{set\_voltage} & \texttt{SOUR:VOLT} & Programa a tensão \\
\texttt{measure\_current} & \texttt{MEAS:CURR?} & Lê a corrente entregue \\
\texttt{turn\_off\_safely} & \texttt{SOUR:VOLT 0} $\to$ \texttt{OUTP OFF} & Zera e desliga \\
\bottomrule
\end{tabular}
\end{center}

\texttt{turn\_off\_safely} zera a tensão \emph{antes} de desligar a saída, e
engole exceções: é chamado em caminhos de erro, onde falhar ruidosamente
deixaria a fonte energizada.

\subsection{O driver do multímetro}

\classe{DMM4050\_Driver} é mais sensível à ordem das operações. Na abertura:

\begin{lstlisting}[language=Python]
if "SOCKET" in resource_string.upper():
    self.inst.write_termination = '\n'   # conexões por rede exigem
    self.inst.read_termination  = '\n'   # terminação explícita
self.inst.write("SYST:REM")  # 1. força modo remoto
self.inst.write("*CLS")      # 2. limpa o buffer de erros
\end{lstlisting}

Na configuração da medida:

\begin{lstlisting}[language=Python]
self.inst.write("CONF:CURR:DC")
self.inst.write("SENS:CURR:DC:RANG 1e-4")  # faixa FIXA de 100 uA
self.inst.write(f"SENS:CURR:DC:NPLC {nplc}")
\end{lstlisting}

\begin{nota}
A faixa fixa de 100~$\mu$A é essencial. Com auto-range, o instrumento escolhe
uma escala maior e os nanoampères da fotocorrente desaparecem na resolução.
O \texttt{NPLC 10} integra cada leitura por dez ciclos da rede, o que
filtra ruído de 50/60~Hz ao custo de tempo de aquisição.
\end{nota}

Ao fechar, o driver envia \texttt{SYST:LOC} antes de encerrar a comunicação,
devolvendo o controle ao painel frontal do aparelho.

\subsection{Threads de infraestrutura}

Duas operações de rede não podem bloquear a interface, e por isso vivem em
threads próprias:

\begin{description}[style=nextline,leftmargin=0.6cm]
  \item[\classe{ScannerThread}]
    Lista os recursos VISA disponíveis, filtra a fonte pelo identificador de
    fabricante no endereço USB, e testa o endereço de rede do multímetro com um
    \texttt{*IDN?} de timeout curto. Emite duas listas de pares
    (nome amigável, endereço).

  \item[\classe{ValidatorThread}]
    Abre um recurso, pergunta \texttt{*IDN?} e emite o resultado. Endereços que
    começam com \texttt{DEMO::} respondem sem passar por VISA.
\end{description}

\classe{HardwareManager} é o \classe{QObject} que orquestra as duas e
retransmite seus sinais para a interface. Ele mantém referências às threads em
execução --- sem isso, o coletor de lixo do Python destruiria um
\classe{QThread} ainda rodando --- e as solta quando cada thread termina,
usando o sinal \texttt{finished} e um \texttt{wait()} de confirmação. Uma
varredura em curso bloqueia novas até terminar.

\subsection{Seleção de drivers e modo demonstração}

Um único ponto decide qual par de drivers usar:

\begin{lstlisting}[language=Python]
def obter_drivers(modo_demonstracao: bool) -> tuple:
    if modo_demonstracao:
        return MockPWS4323_Driver, MockDMM4050_Driver
    return PWS4323_Driver, DMM4050_Driver
\end{lstlisting}

Como os mocks expõem exatamente a mesma interface, quem consome não precisa
saber qual recebeu: instancia e usa. O \classe{ExperimentWorker} só faz uma
distinção adicional --- não abre o \classe{ResourceManager} do PyVISA quando
está em modo demonstração.

% =============================================================
\section{O modelo físico}

Todo o cálculo vive em \arqv{utils/math\_models.py}, em funções puras que
recebem e devolvem vetores \texttt{numpy}. As constantes usadas são as da
CODATA:

\begin{center}
\begin{tabular}{@{}lll@{}}
\toprule
\textbf{Símbolo} & \textbf{Valor} & \textbf{Grandeza} \\
\midrule
$c$ & $299\,792\,458$~m/s & Velocidade da luz \\
$k_B$ & $1{,}380649 \times 10^{-23}$~J/K & Constante de Boltzmann \\
$h_{\text{ref}}$ & $6{,}62607015 \times 10^{-34}$~J$\cdot$s & Constante de Planck \\
\bottomrule
\end{tabular}
\end{center}

\subsection{Da resistência à temperatura}

A resistência do tungstênio varia com a temperatura segundo um polinômio de
segundo grau, com $T_c$ em graus Celsius:

\begin{equation}
R(T_c) = R_0\,\bigl(1 + \alpha T_c + \beta T_c^2\bigr)
\label{eq:rt}
\end{equation}

\subsubsection*{$R_0$ é a resistência a 0 °C, não a medida a frio}

A equação~\eqref{eq:rt} define $R_0$ como a resistência em $T_c = 0$~°C. O
operador, porém, mede o filamento frio na temperatura \emph{ambiente}, onde a
resistência já é cerca de 13\% maior. Usar a medida diretamente como $R_0$
enviesa todas as temperaturas --- e portanto $h$.

A conversão é a própria equação~\eqref{eq:rt} invertida na temperatura em que
a medida foi feita:

\begin{equation}
R_0 = \frac{R(T_{\text{amb}})}{1 + \alpha T_{\text{amb}} + \beta T_{\text{amb}}^2}
\label{eq:r0}
\end{equation}

É o que faz \texttt{corrigir\_r0\_para\_zero\_celsius}. A interface pede os
dois valores --- a resistência medida e a temperatura em que foi medida --- e
exibe ao vivo o $R_0$ resultante, que é o valor efetivamente usado no cálculo.
Para 1{,}2~$\Omega$ medidos a 25~°C, $R_0 = 1{,}0608~\Omega$.

\subsubsection*{Invertendo o modelo}

O experimento mede $R$ e precisa de $T$. Reorganizando a
equação~\eqref{eq:rt} como uma equação do segundo grau em $T_c$:

\begin{equation}
\underbrace{R_0\beta}_{a}\,T_c^2 \;+\; \underbrace{R_0\alpha}_{b}\,T_c
\;+\; \underbrace{(R_0 - R)}_{c} \;=\; 0
\end{equation}

e resolvendo por Bhaskara, tomando a raiz positiva (a única com sentido
físico, pois $R$ cresce com $T$):

\begin{equation}
T_c = \frac{-R_0\alpha + \sqrt{(R_0\alpha)^2 - 4 R_0\beta\,(R_0 - R)}}{2 R_0 \beta},
\qquad
T[\mathrm{K}] = T_c + 273{,}15
\label{eq:bhaskara}
\end{equation}

É exatamente isso que \texttt{calculate\_temperature} implementa, de forma
vetorizada.

\begin{nota}
\textbf{Onde a inversão engana.} Quando a resistência medida fica
\emph{abaixo} de $R_0$ --- o que acontece em tensões próximas de zero --- a
equação~\eqref{eq:bhaskara} ainda tem solução matemática, mas ela não tem
sentido físico. É a origem das temperaturas da ordem de 77~K que aparecem no
início das varreduras completas. Em casos extremos ($R$ negativa) chega a
produzir kelvin negativo.

O discriminante só fica negativo para resistências muito abaixo de $R_0$;
nesse caso a função devolve NaN, ponto a ponto, em vez de lixo. Mas o caso
comum --- a raiz real e absurda --- não é barrado aqui: quem o descarta é a
seleção de pontos da regressão (seção~\ref{sec:limiar}). A separação é
proposital: esta função inverte o modelo, aquela decide o que é utilizável.
\end{nota}

\subsection{Da temperatura à constante de Planck}

A distribuição de Planck para a radiância espectral de um corpo negro é:

\begin{equation}
I(\lambda,T) = \frac{2\pi c^2 h}{\lambda^5}\,
\frac{1}{e^{hc/\lambda k_B T} - 1}
\end{equation}

Quando $hc/(\lambda k_B T) \gg 1$, o termo $-1$ do denominador é desprezível e
vale a \textbf{aproximação de Wien}:

\begin{equation}
I(\lambda,T) \;\approx\; \frac{2\pi c^2 h}{\lambda^5}\,
\exp\!\left(-\frac{hc}{\lambda k_B T}\right)
\label{eq:wien}
\end{equation}

Nas condições deste experimento ($\lambda = 590$~nm, $T$ até 3000~K) o expoente
vale cerca de 8, e o erro cometido pela aproximação é da ordem de
$e^{-8} \approx 0{,}03\%$ --- desprezível diante das demais fontes de erro.

O LED responde linearmente à radiação que recebe na sua faixa espectral, então
a fotocorrente medida é proporcional a $I(\lambda,T)$. Tomando o logaritmo da
equação~\eqref{eq:wien}:

\begin{equation}
\ln I \;=\; \underbrace{\ln\!\left(\frac{2\pi c^2 h}{\lambda^5}\right)}_{\text{constante}}
\;-\; \frac{hc}{\lambda k_B}\cdot\frac{1}{T}
\end{equation}

Ou seja: \textbf{$\ln I$ contra $1/T$ é uma reta}, cuja inclinação é

\begin{equation}
m = -\frac{hc}{\lambda k_B}
\qquad\Longrightarrow\qquad
\boxed{\;h = -\,\frac{m\,\lambda\,k_B}{c}\;}
\label{eq:h}
\end{equation}

Toda a montagem experimental existe para medir essa inclinação.

\subsection{A regressão}

\texttt{calculate\_planck\_constant} executa a cadeia completa:

\begin{enumerate}
  \item \textbf{Filtra os pontos utilizáveis.} Descarta temperaturas nulas ou
        não finitas (que estourariam o cálculo de $1/T$) e fotocorrentes abaixo
        de um limiar de confiança.
  \item \textbf{Lineariza.} Constrói $x = 1/T$ e $y = \ln I$.
  \item \textbf{Ajusta por mínimos quadrados.} Resolve $y = mx + c$ com
        \texttt{numpy.linalg.lstsq}. Exige ao menos dois pontos válidos; abaixo
        disso devolve zeros.
  \item \textbf{Avalia a qualidade.} Calcula
        $R^2 = 1 - \mathrm{SS}_{\text{res}}/\mathrm{SS}_{\text{tot}}$, com
        guarda para o caso $\mathrm{SS}_{\text{tot}} = 0$ (todos os $\ln I$
        iguais), em que $R^2$ é indefinido e se reporta zero.
  \item \textbf{Extrai $h$} pela equação~\eqref{eq:h} e calcula o erro relativo
        contra $h_{\text{ref}}$.
\end{enumerate}

Retorna a tupla \texttt{(h, erro\_relativo, m, c, r2)}.

\subsection{O limiar de confiança e a região de Wien}
\label{sec:limiar}

O passo 1 é o mais delicado do software, e merece detalhamento.

A fotocorrente cai exponencialmente com a queda da temperatura. Abaixo de
aproximadamente 1800~K ela se torna menor que o ruído de fundo do multímetro:
o que se lê não é mais sinal, é o piso do instrumento. Esses pontos não carregam
informação sobre $h$, mas entram na regressão como se carregassem, e a
distorcem.

A seleção é feita por uma função dedicada, \texttt{selecionar\_pontos\_validos},
que aplica três critérios e devolve uma máscara booleana:

\begin{enumerate}[nosep]
  \item \textbf{Temperatura fisicamente possível} --- descarta NaN e qualquer
        $T \le 0$;
  \item \textbf{Região de Wien} --- descarta $T$ abaixo de um limite
        configurável, por padrão 1800~K;
  \item \textbf{Fotocorrente acima do limiar} --- mantém o logaritmo finito.
\end{enumerate}

A mesma função alimenta a interface, que pinta de cinza os pontos descartados:
o operador vê, ponto a ponto, o que entrou e o que não entrou na conta.

\begin{lstlisting}[language=Python]
temperatura_utilizavel = np.isfinite(T_kelvin) & (T_kelvin > 0)
regiao_de_wien        = T_kelvin >= t_minima
sinal_acima_do_piso   = np.isfinite(photocurrent) & (photocurrent > limiar_corrente)
return temperatura_utilizavel & regiao_de_wien & sinal_acima_do_piso
\end{lstlisting}

O efeito do critério pode ser medido. Numa varredura completa de 0 a 12~V com
49 pontos, gerada pela bancada simulada:

\begin{center}
\begin{tabular}{@{}lrrrr@{}}
\toprule
\textbf{Critério de corte} & \textbf{Pontos} & \textbf{$h$ (J$\cdot$s)} & \textbf{Erro} & \textbf{$R^2$} \\
\midrule
Nenhum & 49 & $6{,}9 \times 10^{-36}$ & 99{,}0\% & 0{,}062 \\
Só $I > 1$~nA & 49 & $6{,}9 \times 10^{-36}$ & 99{,}0\% & 0{,}062 \\
$I > 25$~nA & 27 & $6{,}32 \times 10^{-34}$ & 4{,}6\% & 0{,}999 \\
$T > 1800$~K \textbf{(padrão)} & 29 & $6{,}31 \times 10^{-34}$ & 4{,}8\% & 0{,}999 \\
$T > 2000$~K & 23 & $6{,}42 \times 10^{-34}$ & 3{,}1\% & 0{,}999 \\
\bottomrule
\end{tabular}
\end{center}

\begin{nota}
\textbf{Leitura do quadro.} Um limiar de 1~nA sobre a corrente \emph{não
descarta nenhum ponto}: o próprio piso do multímetro é da ordem de 4~nA,
portanto maior que o limiar. É por isso que o corte útil é o de
\emph{temperatura}, e não o de corrente. Um limiar de corrente coerente com o
instrumento precisaria vir do modelo de erro dele --- a especificação de
exatidão na faixa de 100~$\mu$A tem um termo de fundo de 25~nA --- e não de um
número escolhido à mão; essa derivação é trabalho ainda planejado.
\end{nota}

\begin{nota}
\textbf{O ponto central de operação.} O corte acontece na \emph{regressão},
não na \emph{coleta}. A varredura pode cobrir a faixa inteira, inclusive a
partir de 0~V, e o CSV guarda todos os pontos; apenas a conta final usa a
região informativa. Não é preciso escolher entre registrar o comportamento
completo do filamento e obter um bom valor de $h$.
\end{nota}

\subsection{A simulação analítica}

\texttt{simulate\_experiment\_data} gera dados sintéticos para a aba de
Simulação. Ela impõe um vetor de temperaturas linearmente espaçado entre 1500 e
3000~K, calcula a resistência correspondente pela equação~\eqref{eq:rt}, obtém
a corrente por $i = V/R$, e aplica a equação~\eqref{eq:wien} com $h_{\text{ref}}$
para produzir a fotocorrente. Sobre o resultado soma ruído gaussiano
proporcional ao sinal mais um piso instrumental, e limita os valores por baixo
para não quebrar o logaritmo.

\begin{nota}
\textbf{Limite conhecido.} A temperatura é imposta \emph{independentemente} das
tensões: a simulação não reproduz a relação física entre tensão aplicada e
temperatura atingida. Isso é adequado para demonstrar a cadeia de cálculo, mas
não para estudar o comportamento do filamento. Para esse fim existe a bancada
simulada (seção~\ref{sec:mock}), que resolve o balanço de energia de verdade.
\end{nota}

% =============================================================
\section{A bancada simulada}
\label{sec:mock}

\arqv{core/mock\_hardware.py} implementa um filamento e um LED virtuais que
permitem exercitar todo o software sem instrumentos.

\subsection{Por que os dois drivers compartilham estado}

Na bancada real, fonte e multímetro estão acoplados pelo filamento: a tensão
que a fonte aplica define a temperatura, e é a temperatura que o LED enxerga.
Dois mocks independentes não conseguiriam reproduzir isso. Por isso ambos
apontam para uma instância única de \classe{BancadaSimulada}, que guarda o
estado físico.

\subsection{O modelo térmico}

Em regime estacionário, a potência elétrica entregue ao filamento iguala a
potência dissipada:

\begin{equation}
\frac{V^2}{R(T)} \;=\; k_{\text{rad}}\bigl(T^4 - T_{\text{amb}}^4\bigr)
\;+\; k_{\text{cond}}\bigl(T - T_{\text{amb}}\bigr)
\label{eq:balanco}
\end{equation}

O primeiro termo é a radiação (lei de Stefan--Boltzmann); o segundo agrupa as
perdas por condução e convecção, que dominam em temperaturas baixas. Com
$R(T)$ dado pela equação~\eqref{eq:rt}, a diferença entre os dois lados é
monótona decrescente em $T$, o que permite resolver por bissecção sem precisar
de derivada:

\begin{lstlisting}[language=Python]
def desequilibrio(t):
    return tensao**2 / self.resistencia(t) - (
        self.k_rad * (t**4 - T_AMBIENTE**4)
        + self.k_cond * (t - T_AMBIENTE))
\end{lstlisting}

\subsection{Como as constantes foram obtidas}

$k_{\text{rad}}$ e $k_{\text{cond}}$ não foram arbitrados. Foram ajustados por
mínimos quadrados sobre os 1127 pontos úteis dos 52 arquivos CSV de coletas
reais guardados em \texttt{data\_backup/}, resultando em resíduo mediano de
2{,}9\% na potência. A 2500~K, o termo radiativo responde por 99\% da
dissipação --- coerente com a física esperada.

Do ajuste preserva-se a \emph{razão} $k_{\text{cond}}/k_{\text{rad}}$, que
descreve a forma da curva. A escala é recalibrada na construção do objeto, por
bissecção, para que o filamento virtual atinja 2540~K em 12~V --- o topo
efetivamente medido na bancada real.

A constante do LED é fixada de modo que a fotocorrente na âncora valha
$7{,}2 \times 10^{-7}$~A, o valor máximo observado nas coletas reais. O
expoente usa $h_{\text{ref}}$: os dados gerados \emph{contêm} o valor
verdadeiro de $h$, e recuperá-lo é o teste de que a cadeia de cálculo está
correta.

\subsection{Erros e ruídos reproduzidos}

\begin{center}
\begin{longtable}{@{}p{5.0cm}p{9.6cm}@{}}
\toprule
\textbf{Efeito} & \textbf{Como é modelado} \\
\midrule
\endhead
Erro de programação da fonte & Distribuição retangular sobre o valor pedido;
o valor efetivamente aplicado difere do programado. \\[0.3em]
Erro de leitura da fonte & Distribuição retangular, com termos proporcional e
fixo, separadamente para tensão e corrente. \\[0.3em]
Resolução do multímetro & Leitura arredondada para passos de 100~pA. \\[0.3em]
Piso do multímetro & Um deslocamento sistemático de alguns nanoampères mais
ruído gaussiano --- reproduz o que os CSVs reais mostram em tensão baixa. \\[0.3em]
Modo corrente constante & Se a corrente calculada excede o limite configurado,
a fonte segura a corrente e a tensão nos terminais cai. \\[0.3em]
Saída desligada & Corrente nula e temperatura de volta ao ambiente. \\
\bottomrule
\end{longtable}
\end{center}

O gerador aleatório aceita semente, o que torna uma varredura inteiramente
reprodutível --- útil para testes de regressão.

\begin{nota}
\textbf{O piso é reproduzido de propósito.} Sem ele, o mock esconderia
justamente o problema descrito na seção~\ref{sec:limiar}, e não haveria como
verificar a correção quando ela for implementada.
\end{nota}

\subsection{O que o mock ainda não reproduz}

Inércia térmica --- o filamento virtual salta direto ao regime estacionário,
o que é aceitável porque os workers já esperam o tempo de estabilização --- e a
resistência dos cabos da medição a dois fios.

% =============================================================
\section{Concorrência}

\subsection{Por que threads}

Uma coleta real leva minutos e é dominada por esperas: estabilização térmica,
ida e volta de cada comando SCPI. Fazer isso na thread da interface congelaria
a janela. Cada coleta roda, portanto, em um \classe{QThread} próprio, que se
comunica com a interface exclusivamente por sinais.

\begin{nota}
Os drivers VISA são instanciados \emph{dentro} da thread que os usa. Isso é
deliberado: conexões VISA não devem ser compartilhadas entre threads.
\end{nota}

\subsection{Os sinais}

\begin{center}
\begin{tabular}{@{}p{3.4cm}p{4.0cm}p{7.0cm}@{}}
\toprule
\textbf{Sinal} & \textbf{Carga} & \textbf{Quando} \\
\midrule
\texttt{new\_data\_point} & tensão, temperatura, fotocorrente & A cada ponto medido \\[0.3em]
\texttt{finished\_exp} & $h$, erro, $m$, $c$, $R^2$ & Ao concluir a regressão \\[0.3em]
\texttt{error\_occurred} & mensagem & Em falha de comunicação ou parada precoce \\
\bottomrule
\end{tabular}
\end{center}

\subsection{Parada segura}

\texttt{stop()} apenas baixa uma flag; a thread a consulta no topo de cada
iteração e sai do laço por conta própria. A interface \emph{não} espera pela
thread nesse momento --- ela continua respondendo, e a thread desliga a fonte,
calcula o resultado parcial e emite \texttt{finished\_exp} naturalmente.

O bloco \texttt{finally} do laço de coleta é o que garante a segurança da
bancada: aconteça o que acontecer, a fonte é zerada, desligada e fechada, e o
multímetro devolvido ao modo local.

\begin{lstlisting}[language=Python]
finally:
    if pws:
        pws.turn_off_safely()   # zera a tensão e desliga a saída
        pws.close()
    if dmm:
        dmm.close()             # devolve o painel frontal ao operador
\end{lstlisting}

% =============================================================
\section{O fluxo de um experimento}

Passo a passo, do clique ao resultado:

\begin{enumerate}
  \item A interface lê os nove campos de parâmetros, converte para número e
        monta um dicionário.
  \item Lê do \classe{QSettings} o modo demonstração e os endereços dos
        instrumentos.
  \item Dimensiona a barra de progresso pelo \emph{mesmo} vetor de tensões que
        o worker vai percorrer --- calcular o total por outra fórmula produz
        divergências por aritmética de ponto flutuante.
  \item Cria o \classe{ExperimentWorker}, conecta os três sinais e inicia a
        thread.
  \item O worker escolhe o par de drivers, instancia ambos, configura o limite
        de corrente da fonte e a faixa/integração do multímetro.
  \item Escreve o cabeçalho do CSV.
  \item Liga a saída da fonte e entra no laço. Para cada tensão:
    \begin{enumerate}[label=\alph*.]
      \item programa a tensão;
      \item aguarda o tempo de estabilização térmica;
      \item lê a corrente na fonte e a fotocorrente no multímetro;
      \item calcula $R = V/i$, com piso na corrente para evitar divisão por zero;
      \item converte $R$ em $T$ pela equação~\eqref{eq:bhaskara};
      \item grava a linha no CSV \emph{imediatamente}, com abertura e
            fechamento do arquivo a cada ponto;
      \item emite \texttt{new\_data\_point}, e a interface atualiza gráficos,
            registro e barra.
    \end{enumerate}
  \item Ao fim (ou na parada), desliga a saída e, havendo mais de dois pontos,
        chama \texttt{calculate\_planck\_constant} e emite o resultado.
  \item O bloco \texttt{finally} garante o desligamento seguro.
\end{enumerate}

\begin{nota}
\textbf{Sobre a resistência calculada.} A resistência do filamento vem da
tensão \emph{lida de volta} da fonte (\texttt{MEAS:VOLT?}), não do valor
programado: o readback do PWS4323 é mais exato que a exatidão de programação.
Se o instrumento não responder ao readback, o cálculo cai para o valor
programado em vez de abortar a coleta.

Como a medição é a dois fios, a resistência dos cabos entra somada à do
filamento; a interface oferece um campo para descontá-la:
\[
R_{\text{filamento}} = \frac{V_{\text{medida}}}{i} - R_{\text{cabos}}
\]
As duas tensões --- programada e medida --- são gravadas no CSV, o que permite
auditar depois a diferença entre elas.
\end{nota}

% =============================================================
\section{Formato dos dados}

\subsection{O CSV de coleta}

Um arquivo por execução, em \texttt{Software/data\_backup/}:

\begin{center}
\begin{tabular}{@{}ll@{}}
\toprule
Coletas reais & \texttt{exp\_planck\_AAAAMMDD\_HHMMSS.csv} \\
Coletas simuladas & \texttt{demo\_planck\_AAAAMMDD\_HHMMSS.csv} \\
\bottomrule
\end{tabular}
\end{center}

Prefixos distintos impedem que dados de demonstração contaminem o acervo de
medidas reais --- que é, entre outras coisas, a base de calibração da bancada
simulada.

\begin{center}
\begin{tabular}{@{}lll@{}}
\toprule
\textbf{Coluna} & \textbf{Unidade} & \textbf{Origem} \\
\midrule
\texttt{Tensao\_Fonte\_V} & V & Valor programado na fonte \\
\texttt{Corrente\_Filamento\_A} & A & \texttt{MEAS:CURR?} na fonte \\
\texttt{Resistencia\_Ohms} & $\Omega$ & Calculada, $V_{\text{medida}}/i - R_{\text{cabos}}$ \\
\texttt{Temperatura\_K} & K & Calculada, equação~\eqref{eq:bhaskara} \\
\texttt{Fotocorrente\_A} & A & \texttt{READ?} no multímetro \\
\texttt{Tensao\_Medida\_V} & V & \texttt{MEAS:VOLT?} na fonte \\
\bottomrule
\end{tabular}
\end{center}

\begin{nota}
Há dados históricos gravados neste formato. A coluna de tensão medida foi
acrescentada \emph{ao final}, justamente para que as cinco colunas originais
permaneçam na mesma ordem: arquivos antigos continuam legíveis, e leitores que
acessam por posição não quebram. Qualquer alteração futura de esquema deve
seguir a mesma regra.
\end{nota}

\subsection{O relatório PDF}

\texttt{generate\_planck\_report} monta o documento com \texttt{reportlab},
recebendo quatro blocos: resumo de resultados, dicionário de parâmetros,
metadados informados pelo usuário e o caminho de uma imagem. A imagem é uma
captura da área de gráficos, gerada no momento da exportação e apagada em
seguida.

% =============================================================
\section{Camada de conteúdo}

\arqv{content/referencias.py} guarda a lista de referências como dados tipados,
não como código de interface. Cada item é uma \texttt{dataclass}
\classe{Referencia} congelada, com título, autoria, publicação, identificador,
nome do arquivo local, URLs oficiais, condição de acesso, o uso que o software
faz do documento e a lista de trechos aproveitados.

Duas propriedades calculadas resolvem o caminho do arquivo na pasta
\texttt{Docs/} e informam se ele está presente --- a interface usa isso para
desabilitar o botão de cópia local quando o PDF não existe.

Para acrescentar um documento, basta acrescentar um item à lista: a aba se
adapta sozinha.

% =============================================================
\section{Aparência}

\arqv{ui/theme.py} exporta uma única string QSS aplicada globalmente. A paleta é
escura, com fundo quase preto, superfícies em cinza-escuro e azul como cor de
destaque. Os gráficos do \texttt{pyqtgraph} são configurados separadamente, com
fundo e primeiro plano combinando com o tema.

\begin{nota}
Alguns componentes ainda trazem estilo embutido no próprio código, em vez de
virem da folha global. É um ponto de dívida técnica conhecido, a resolver numa
revisão de interface.
\end{nota}

% =============================================================
\section{Como estender}

\begin{description}[style=nextline,leftmargin=0.6cm]

  \item[Acrescentar uma referência]
    Adicione um item à lista em \arqv{content/referencias.py}, com o nome do
    arquivo dentro de \texttt{Docs/}. Nada mais precisa mudar.

  \item[Suportar outro instrumento]
    Escreva uma classe com os mesmos métodos públicos do driver existente e
    faça \texttt{obter\_drivers} devolvê-la. Nenhum outro arquivo precisa saber.
    O arquivo de testes verifica automaticamente que a interface do mock cobre
    a do driver real --- o mesmo critério vale para um driver novo.

  \item[Alterar o modelo físico]
    Todas as contas estão em \arqv{utils/math\_models.py}, em funções puras sem
    dependência de Qt. São testáveis isoladamente.

  \item[Acrescentar uma aba]
    Crie o widget em \arqv{ui/tabs/} e registre-o em \texttt{setup\_ui}, na
    janela principal.

\end{description}

% =============================================================
\section{Testes}

\texttt{Tests/test\_mock\_hardware.py} roda sem hardware e sem dependências
extras:

\begin{quote}\ttfamily
cd Software \&\& python ../Tests/test\_mock\_hardware.py
\end{quote}

Cobre quatro grupos de verificação: que os mocks são substituíveis pelos
drivers reais (comparando as interfaces por introspecção), que a física
simulada está calibrada nos pontos de âncora e é monótona, que os
comportamentos do instrumento aparecem (piso do multímetro, modo corrente
constante, corte seguro), e que a cadeia de cálculo recupera $h$ dos dados
gerados.

Um dos testes fixa deliberadamente o \emph{erro esperado} numa varredura
completa. Quando o critério de corte descrito na seção~\ref{sec:limiar} for
corrigido, esse teste falhará --- e é assim que se saberá que a correção
funcionou.

\vfill
\begin{center}
\rule{0.5\linewidth}{0.4pt}\\[0.5em]
\small Documento companheiro: \textbf{Manual do Usuário}, com o procedimento de
operação.
\end{center}

\end{document}
